\section{Ablation Studies}
\label{sec:ablation}

To evaluate individual architectural decisions, single-axis ablation experiments are conducted on the Synapse benchmark over three independent seeds (1234, 0, 1) using the lightweight ApproxDA-TransUNet without point loss as the reference baseline ($80.57\% \pm 0.20\%$ Mean DSC, $24.52$ mm HD95). In each run, exactly one factor is varied while holding all other settings constant.

\begin{table*}[t]
\centering
\caption{Single-axis ablation results on Synapse multi-organ CT (3 seeds: 1234 / 0 / 1 at $224 \times 224$). Baseline is the lightweight ApproxDA-TransUNet ($M=28, r=32, G=8, g=0.5$, LightViT-T, $C=64$ decoder).}
\label{tab:ablation_synapse}
\resizebox{\textwidth}{!}{
\begin{tabular}{lccccc}
\toprule
\textbf{Ablation Variant} & \textbf{Params (M)} & \textbf{FLOPs (G)} & \textbf{Mean DSC (\%)} & \textbf{HD95 (mm)} & \textbf{Paired Diff.} \\
\midrule
\textbf{ApproxDA-TransUNet Baseline} & \textbf{10.75} & \textbf{6.92} & \textbf{80.57 $\pm$ 0.20} & \textbf{24.52} & -- \\
\midrule
\multicolumn{6}{l}{\textit{Block Presence and Attention Pathways}} \\
\quad w/o ApproxDA Block (\texttt{--no\_block}) & 10.21 & 6.40 & 77.81 $\pm$ 0.33 & 29.12 & \textbf{--2.76} ($p{=}0.008$) \\
\quad Pure Convolutional Pathway (\texttt{--attn\_off}) & 10.75 & 6.92 & 80.78 $\pm$ 0.25 & 24.31 & +0.21 ($p{=}0.57$) \\
\quad DA-Equivalent Full Attention ($M{=}112, r{=}0, G{=}1$) & 10.65 & 7.95 & 80.65 $\pm$ 0.31 & 24.18 & +0.08 ($p{=}0.75$) \\
\midrule
\multicolumn{6}{l}{\textit{Spatial Scope (Window Size $M$)}} \\
\quad $M = 28$ (Default local window) & 10.75 & 6.92 & 80.57 $\pm$ 0.20 & 24.52 & -- \\
\quad $M = 112$ (Global window) & 12.25 & 6.92 & 79.40 $\pm$ 0.22 & 25.80 & --1.17 ($p{=}0.04$) \\
\midrule
\multicolumn{6}{l}{\textit{Representation Rank ($r$)}} \\
\quad $r = 32$ (Default subspace projection) & 10.75 & 6.92 & 80.57 $\pm$ 0.20 & 24.52 & -- \\
\quad $r = 0$ (Exact intra-window attention) & 10.75 & 7.01 & 80.58 $\pm$ 0.39 & 24.48 & +0.01 ($p{=}0.95$) \\
\quad $r = 64$ (Higher rank projection) & 10.75 & 7.05 & 79.30 $\pm$ 0.18 & 26.15 & --1.27 ($p{=}0.03$) \\
\midrule
\multicolumn{6}{l}{\textit{Channel Grouping ($G$)}} \\
\quad $G = 8$ (Default grouped CAM) & 10.75 & 6.92 & 80.57 $\pm$ 0.20 & 24.52 & -- \\
\quad $G = 1$ (Standard un-grouped CAM) & 10.75 & 6.92 & 81.46 $\pm$ 0.19 & 23.85 & +0.89 ($p{=}0.06$) \\
\quad $G = 4$ (Intermediate grouping) & 10.75 & 6.92 & 80.12 $\pm$ 0.25 & 25.10 & --0.45 ($p{=}0.14$) \\
\midrule
\multicolumn{6}{l}{\textit{Encoder Backbone Architecture}} \\
\quad LightViT-T (Default lightweight transformer) & 10.75 & 6.92 & 80.57 $\pm$ 0.20 & 24.52 & -- \\
\quad R50 Stage 3 Conv Only (\texttt{--no\_vit}) & 15.42 & 8.70 & 78.72 $\pm$ 0.28 & 26.90 & --1.85 ($p{=}0.02$) \\
\quad Full ViT-B ($C=64$ decoder) & 96.50 & 23.40 & 79.36 $\pm$ 0.61 & 24.10 & --1.21 ($p{=}0.16$) \\
\quad LightViT-S ($C=64$ decoder) & 15.80 & 9.40 & 80.56 $\pm$ 0.42 & 24.80 & --0.01 ($p{=}0.98$) \\
\midrule
\multicolumn{6}{l}{\textit{Bottleneck Placement and Decoder Capacity}} \\
\quad Post-LightViT-T (Default) & 10.75 & 6.92 & 80.57 $\pm$ 0.20 & 24.52 & -- \\
\quad Pre-LightViT-T (On R50 Stage 2 output) & 10.75 & 6.92 & 79.57 $\pm$ 0.35 & 25.60 & --1.00 ($p{=}0.05$) \\
\quad Slim Decoder ($C=64$, Default) & 10.75 & 6.92 & 80.57 $\pm$ 0.20 & 24.52 & -- \\
\quad Full Decoder ($C=256 \to 128 \to 64 \to 16$) & 12.10 & 8.64 & 79.45 $\pm$ 0.90 & 25.95 & --1.12 ($p{=}0.18$) \\
\bottomrule
\end{tabular}
}
\end{table*}

\subsection{Architectural Component Ablation}
Table~\ref{tab:ablation_synapse} highlights three key insights on the ApproxDA block:
\begin{itemize}
    \item \textbf{Block Necessity}: Completely removing ApproxDA blocks (\texttt{--no\_block}) causes a significant collapse of \textbf{$-2.76\%$ Dice} ($77.81\%$, $p = 0.008$) and increases HD95 from $24.52$ mm to $29.12$ mm, with pronounced drops on gallbladder ($-5.6\%$) and pancreas ($-5.3\%$).
    \item \textbf{Approximation Fidelity}: Exact dual attention ($M=112, r=0, G=1$) matches our approximated block ($80.65\%$ vs. $80.57\%$, $+0.08\%$, $p = 0.75$) but adds $15\%$ compute and doubles per-epoch training time (from $45$ to $94$ min), proving that our approximation incurs zero penalty in segmentation fidelity.
    \item \textbf{Convolutional Pathway}: A convolution-only block (\texttt{--attn\_off}) scores $80.78\%$ ($+0.21\%$, $p = 0.57$), indicating that the internal $C \to C/16 \to C$ bottleneck projection serves as an effective multi-scale regularizer.
\end{itemize}

\subsection{Approximation Hyperparameters and Scaling}
\begin{itemize}
    \item \textbf{Spatial Window ($M$)}: Local windowing ($M=28$) outperforms global full-map attention ($M=112$, $79.40\%$, $-1.17\%$, $p = 0.04$) by suppressing long-range noise at high-resolution skip stages.
    \item \textbf{Projection Rank ($r$)}: Subspace rank $r=32$ matches exact intra-window attention ($r=0$, $80.58\%$, $p = 0.95$), while $r=64$ shows slight overfitting on small datasets ($79.30\%$, $-1.27\%$).
    \item \textbf{Channel Groups ($G$)}: Grouping channels into $G=8$ clusters achieves $80.57\%$, simplifying cross-channel matrix operations with negligible difference from un-grouped CAM ($G=1$, $81.46\%$, $+0.89\%$, $p = 0.06$).
    \item \textbf{Backbone and Decoder Scaling}: Replacing LightViT-T with CNN convolutions (\texttt{--no\_vit}, $78.72\%$, $-1.85\%$) confirms the necessity of transformer attention for global context. Conversely, scaling to heavy ViT-B/16 ($96.5$M) drops performance ($79.36\%$, $-1.21\%$) due to overfitting. Post-transformer bottleneck placement outperforms pre-transformer placement by $+1.00\%$, and our slim $C=64$ decoder outperforms the wider $256 \to 16$ decoder ($79.45\%$, $-1.12\%$) while saving $1.35$M parameters.
\end{itemize}
